\section{A Working Church: Baptism, Prayer, Penance, Marriage}

Because Tertullian wrote for catechumens, spouses, and penitents as well as
against enemies, his corpus is an unmatched window on what a Latin church
around the year 200 actually did. The evidence must be read as advocacy---he is
always urging a practice, rarely just describing one---but the practices
themselves are not invented.

\subsection{The font}

\work{On Baptism}, the earliest surviving treatise on any sacrament, defends
the humble efficacy of water against a female teacher who denied it.
``We, little fishes, after the example of our \textit{Ichthys} Jesus
Christ, are born in water,'' he begins---the Greek acrostic ``fish'' standing
for Christ---``nor have we safety in any other way than by permanently abiding
in water'' (\work{On Baptism} 1). The treatise walks through a
developed rite. The water itself is no mere symbol: invoked over, it is
changed, for ``the nature of the waters, sanctified by the Holy One, itself
conceived withal the power of sanctifying'' (\work{On Baptism} 4). Then the
sequence that every later Western ritual would elaborate: ``After this, when
we have issued from the font, we are thoroughly anointed with a blessed
unction'' (\work{On Baptism} 7), and ``the hand is laid on us, invoking and
inviting the Holy Spirit through benediction'' (\work{On Baptism} 8)---
washing, anointing, imposition of hands, each already distinct, each already
interpreted. The same treatise preserves, in passing, an episode of church
discipline elsewhere: the presbyter in Asia who composed the apocryphal
\work{Acts of Paul}, ``after being convicted, and confessing that he
had done it from love of Paul, was removed from his office'' (\work{On Baptism}
17), an early datum for how the church policed writings about the apostles.

One counsel in the treatise startles the modern reader and is precious
evidence precisely for that reason: Tertullian argues that ``the delay of
baptism is preferable; principally, however, in the case of little
children\ldots\ Let them become Christians when they have become able to know
Christ'' (\work{On Baptism} 18). The passage presupposes that infants were in
fact being brought to the font---otherwise there would be nothing to
delay---and simultaneously shows that the practice was still arguable; his
motive is not doubt about the sacrament but terror of sin after it, the same
once-only seriousness that governs his penitential teaching. Both halves of
the datum, the practice and the protest, belong to the history of infant
baptism.

\subsection{The school of prayer}

\work{On Prayer} expounds the Our Father for catechumens, clause by clause,
and opens with the characterization the Catholic Church still quotes: the
Lord, he says, ``has determined for us, the disciples of the New Testament, a
new form of prayer,'' and ``in the Prayer is comprised an epitome of the whole
Gospel'' (\work{On Prayer} 1). The treatise then descends into the
congregation's actual habits---postures, the kiss of peace, disputes about
sitting on couches and washing hands---small windows onto a worship already
thick with custom, including the earliest Latin notice of the fixed hours of
daily prayer: beyond morning and evening, ``those common hours, I mean, which
mark the intervals of the day---the third, the sixth, the ninth---which we
may find in the Scriptures to have been more solemn than the rest''
(\work{On Prayer} 25)---the seed of what the later church would organize as
the daily office. It closes on the note that makes it a
devotional classic rather than a rubric-book: ``Prayer is alone that which
vanquishes God'' (\work{On Prayer} 29)---prayer as the one force to which,
by his own covenant, the Almighty consents to yield.

\subsection{The second plank}

\work{On Repentance} shows the second-century penitential system mid-development.
Baptism is the great forgiveness; for grave sin afterwards, God ``in the
vestibule has stationed the second repentance for opening to such as knock: but
now once for all, because now for the second time; but never more'' (\work{On
Repentance} 7). That second chance takes bodily form in \textit{exomologesis},
the public discipline in which the sinner is commanded ``to lie in sackcloth
and ashes, to cover his body in mourning,'' and ``to bow before the feet of the
presbyters'' (\work{On Repentance} 9). The whole later history of the sacrament
of penance---and Tertullian's own later fury at its relaxation---starts from
this once-only public rigor.

\subsection{Custom with the force of law}

\work{The Chaplet} preserves the classic early list of practices held on custom
and tradition rather than on a scriptural text: the solemn renunciation of
``the devil, and his pomp, and his angels'' before baptism; the threefold
immersion (``Hereupon we are thrice immersed, making a somewhat ampler pledge
than the Lord has appointed in the Gospel''); the taste of milk and honey for
the newly baptized; the eucharistic assembly itself---``we take also, in
congregations before daybreak, and from the hand of none but the presidents,
the sacrament of the Eucharist, which the Lord both commanded to be eaten at
meal-times, and enjoined to be taken by all alike''---with its pre-dawn hour
and its rule that only the presiding clergy distribute; and two habits
of enormous later history---``we make offerings for the dead as birthday
honours,'' that is, the Eucharist offered on the anniversary of a death, and
the constant signing of the cross: ``At every forward step and movement, at
every going in and out\ldots\ in all the ordinary actions of daily life, we
trace upon the forehead the sign'' (\work{The Chaplet} 3). Tertullian cites
these precisely because they are unwritten, arguing that observed tradition can
bind; the passage remains a locus classicus in Catholic accounts of tradition.

\subsection{The household}

The household had its theology too. The second book \work{To His Wife} closes
with the most quoted early Christian praise of marriage: ``What kind of yoke is
that of two believers, (partakers) of one hope, one desire, one discipline, one
and the same service?\ldots\ truly two in one flesh. Where the flesh is one,
one is the spirit too''; they pray and fast together, and ``When Christ sees
and hears, He joys'' (\work{To His Wife} 2.8). The same pen that could scorch
could also bless; the blessing, characteristically, is for a marriage inside
the discipline of the church, contracted before it and confirmed by its
oblation. Even the ferocious books on women's dress end not in prohibition but
in a transposed finery: ``Clothe yourselves with the silk of uprightness, the
fine linen of holiness, the purple of modesty. Thus painted, you will have God
as your Lover!'' (\work{On the Apparel of Women} 2.13).

\subsection{The Christian in a pagan city}

Between font and household lay the city, and the city was wired with religion.
\work{On Idolatry} is Tertullian's casuistry manual for living in it, built on
a maximal definition of the danger: ``The principal crime of the human race,
the highest guilt charged upon the world, the whole procuring cause of
judgment, is idolatry'' (\work{On Idolatry} 1). From that premise he works
through the trades one by one---makers of idols, teachers of pagan letters,
astrologers, incense-sellers, holders of civic office---asking of each whether
it can be practiced without serving the gods. He is no separatist in
principle: Scripture, he concedes, does not forbid ordinary commerce and
company with pagans, ``Otherwise you would go out from the world'' (\work{On
Idolatry} 14, citing Paul). But where the line falls, it falls absolutely.
On military service his ruling became famous: soldiers had come to John the
Baptist, a centurion had believed, ``still the Lord afterward, in disarming
Peter,'' unbelted---so the standard rendering---every soldier thereafter
(\work{On Idolatry} 19). He pressed the same conclusion when a Christian
soldier's refusal to wear the military crown scandalized his fellow
believers: ``shall it be held lawful to make an occupation of the sword, when
the Lord proclaims that he who uses the sword shall perish by the sword? And
shall the son of peace take part in the battle when it does not become him
even to sue at law?'' (\work{The Chaplet} 11). Later Christian
centuries would decide the question differently; the sentences remain the
sharpest ancient statement of the rigorist side. \work{The Shows} applies
the same total logic to the games and the theatre, and its legal ground is
the baptismal contract itself: ``When entering the water, we make profession
of the Christian faith in the words of its rule; we bear public testimony
that we have renounced the devil, his pomp, and his angels''---and the
shows, being idolatry's pageantry, are precisely that renounced pomp
(\work{The Shows} 4). The conclusion is drawn with characteristic relish:
``How monstrous it is to go from God's church to the devil's---from the sky
to the stye, as they say; to raise your hands to God, and then to weary them
in the applause of an actor'' (\work{The Shows} 25). Whether all this describes
what average Carthaginian Christians did, or what Tertullian could not make
them do, the treatises themselves confess: he legislates because they
attended, traded, served, and wore what he condemned.
